\chapter{2020年全国II卷（文）}

\section{单选题}

\begin{problem}
已知集合 \(A=\{x\mid \lvert x\rvert<3,\ x\in \Z\}\)，\(B=\{x\mid \lvert x\rvert>1,\ x\in \Z\}\)，则 \(A\cap B=\)（\quad）

\choices
    {\(\varnothing\)}
    {\(\{-3,-2,2,3\}\)}
    {\(\{-2,0,2\}\)}
    {\(\{-2,2\}\)}
\end{problem}

\begin{problem}
\((1-\i)^4=\)（\quad）

\choices
    {\(-4\)}
    {\(4\)}
    {\(-4\i\)}
    {\(4\i\)}
\end{problem}

\begin{problem}
如图，将钢琴上的 \(12\) 个键依次记为 \(a_1\)，\(a_2\)，\(\cdots\)，\(a_{12}\). 设 \(1\le i<j<k\le 12\). 若 \(k-j=3\) 且 \(j-i=4\)，则 \(a_i\)，\(a_j\)，\(a_k\) 为原位大三和弦；若 \(k-j=4\) 且 \(j-i=3\)，则称 \(a_i\)，\(a_j\)，\(a_k\) 为原位小三和弦. 用这 \(12\) 个键可以构成的原位大三和弦与原位小三和弦的个数之和为（\quad）

\FigureLayoutDeclare{img/2020/national_paper_2_liberal/q3_fig1.png}{14.0cm}{17bp 19bp 31bp 22bp}
\bitmapfigure[max width=6.0cm]{img/2020/national_paper_2_liberal/q3_fig1.png}

\choices
    {\(5\)}
    {\(8\)}
    {\(10\)}
    {\(15\)}
\end{problem}

\begin{problem}
在新冠肺炎疫情防控期间，某超市开通网上销售业务，每天能完成 \(1200\) 份订单的配货，由于订单量大幅增加，导致订单积压. 为解决困难，许多志愿者踊跃报名参加配货工作. 已知该超市某日积压 \(500\) 份订单未配货，预计第二天的新订单超过 \(1600\) 份的概率为 \(0.05\). 志愿者每人每天能完成 \(50\) 份订单的配货，为使第二天完成积压订单及当日订单的配货的概率不小于 \(0.95\)，则至少需要志愿者（\quad）

\choices
    {\(10\text{ 名}\)}
    {\(18\text{ 名}\)}
    {\(24\text{ 名}\)}
    {\(32\text{ 名}\)}
\end{problem}

\begin{problem}
已知单位向量 \(\bs a\)，\(\bs b\) 的夹角为 \(60^\circ\)，则在下列向量中，与 \(\bs b\) 垂直的是（\quad）

\choices
    {\(\bs a+2\bs b\)}
    {\(2\bs a+\bs b\)}
    {\(\bs a-2\bs b\)}
    {\(2\bs a-\bs b\)}
\end{problem}

\begin{problem}
记 \(S_n\) 为等比数列 \(\{a_n\}\) 的前 \(n\) 项和. 若 \(a_5-a_3=12\)，\(a_6-a_4=24\)，则 \(\dfrac{S_n}{a_n}=\)（\quad）

\choices
    {\(2^n-1\)}
    {\(2-2^{1-n}\)}
    {\(2-2^{n-1}\)}
    {\(2^{1-n}-1\)}
\end{problem}

\begin{problem}
执行下面的程序框图，若输入的 \(k=0\)，\(a=0\)，则输出的 \(k\) 为（\quad）

\bitmapfigure[width=5.0cm]{img/2020/national_paper_2_liberal/q7_fig1.png}

\choices
    {\(2\)}
    {\(3\)}
    {\(4\)}
    {\(5\)}
\end{problem}

\begin{problem}
若过点 \((2,1)\) 的圆与两坐标轴都相切，则圆心到直线 \(2x-y-3=0\) 的距离为（\quad）

\choices
    {\(\dfrac{\sqrt5}{5}\)}
    {\(\dfrac{2\sqrt5}{5}\)}
    {\(\dfrac{3\sqrt5}{5}\)}
    {\(\dfrac{4\sqrt5}{5}\)}
\end{problem}

\begin{problem}
设 \(O\) 为坐标原点，直线 \(x=a\) 与双曲线 \(C:\dfrac{x^2}{a^2}-\dfrac{y^2}{b^2}=1\)（\(a>0\)，\(\ b>0\)）的两条渐近线分别交于 \(D\)，\(E\) 两点. 若 \(\triangle ODE\) 的面积为 \(8\)，则 \(C\) 的焦距的最小值为（\quad）

\choices
    {\(4\)}
    {\(8\)}
    {\(16\)}
    {\(32\)}
\end{problem}

\begin{problem}
设函数 \(f(x)=x^3-\dfrac{1}{x^3}\)，则 \(f(x)\)（\quad）

\choices
    {是奇函数，且在 \((0,+\infty)\) 单调递增}
    {是奇函数，且在 \((0,+\infty)\) 单调递减}
    {是偶函数，且在 \((0,+\infty)\) 单调递增}
    {是偶函数，且在 \((0,+\infty)\) 单调递减}
\end{problem}

\begin{problem}
已知 \(\triangle ABC\) 是面积为 \(\dfrac{9\sqrt3}{4}\) 的等边三角形，且其顶点都在球 \(O\) 的球面上. 若球 \(O\) 的表面积为 \(16\pi\)，则 \(O\) 到平面 \(ABC\) 的距离为（\quad）

\choices
    {\(\sqrt3\)}
    {\(\dfrac{3}{2}\)}
    {\(1\)}
    {\(\dfrac{\sqrt3}{2}\)}
\end{problem}

\begin{problem}
若 \(2^x-2^y<3^{-x}-3^{-y}\)，则（\quad）

\choices
    {\(\ln(y-x+1)>0\)}
    {\(\ln(y-x+1)<0\)}
    {\(\ln\lvert x-y\rvert>0\)}
    {\(\ln\lvert x-y\rvert<0\)}
\end{problem}

\section{填空题}

\begin{problem}
若 \(\sin x=-\dfrac{2}{3}\)，则 \(\cos 2x=\fillinblank{}\).
\end{problem}

\begin{problem}
记 \(S_n\) 为等差数列 \(\{a_n\}\) 的前 \(n\) 项和. 若 \(a_1=-2\)，\(a_2+a_6=2\)，则 \(S_{10}=\fillinblank{}\).
\end{problem}

\begin{problem}
若 \(x\)，\(y\) 满足约束条件
\(\begin{cases}
x+y\ge -1,\\
x-y\ge -1,\\
2x-y\le 1
\end{cases}\)，则 \(z=x+2y\) 的最大值是\fillinblank{}.
\end{problem}

\begin{problem}
设有下列四个命题：

\begin{circlelist}
\item \(p_1\)：两两相交且不过同一点的三条直线必在同一平面内；
\item \(p_2\)：过空间中任意三点有且仅有一个平面；
\item \(p_3\)：若空间两条直线不相交，则这两条直线平行；
\item \(p_4\)：若直线 \(l\subset \alpha\)，直线 \(m\perp \alpha\)，则 \(m\perp l\).
\end{circlelist}

则下述命题中所有真命题的序号是\fillinblank{}

\begin{circlelist}
    \item \(p_1\land p_4\)
    \item \(p_1\land p_2\)
    \item \(\neg p_2\vee p_3\)
    \item \(\neg p_3\vee \neg p_4\)
\end{circlelist}
\end{problem}

\section{解答题}

\begin{problem}
\(\triangle ABC\) 的内角 \(A\)，\(B\)，\(C\) 的对边分别为 \(a\)，\(b\)，\(c\)，已知 \(\cos^2\left( \dfrac{\pi}{2}+A \right)+\cos A=\dfrac{5}{4}\).

\begin{enumerate}
\item 求 \(A\)；
\item 若 \(b-c=\dfrac{\sqrt3}{3}a\)，证明：\(\triangle ABC\) 是直角三角形.
\end{enumerate}
\end{problem}

\begin{problem}
某沙漠地区经过治理，生态系统得到很大改善，野生动物数量有所增加. 为调查该地区某种野生动物的数量，将其分成面积相近的 \(200\) 个地块，从这些地块中用简单随机抽样的方法抽取 \(20\) 个作为样区，调查得到样本数据 \((x_i,y_i)\)（\(i=1\)，\(2\)，\(\cdots\)，\(20\)），其中 \(x_i\) 和 \(y_i\) 分别表示第 \(i\) 个样区的植物覆盖面积（单位：公顷）和这种野生动物的数量，并计算得 \(\sum_{i=1}^{20}x_i=60\)，\(\sum_{i=1}^{20}y_i=1200\)，\(\sum_{i=1}^{20}(x_i-\bar x)^2=80\)，\(\sum_{i=1}^{20}(y_i-\bar y)^2=9000\)，\(\sum_{i=1}^{20}(x_i-\bar x)(y_i-\bar y)=800\).

\begin{enumerate}
\item 求该地区这种野生动物数量的估计值（这种野生动物数量的估计值等于样区这种野生动物数量的平均数乘以地块数）；
\item 求样本 \((x_i,y_i)\)（\(i=1\)，\(2\)，\(\cdots\)，\(20\)）的相关系数（精确到 \(0.01\)）；
\item 根据现有统计资料，各地块间植物覆盖面积差异很大. 为提高样本的代表性以获得该地区这种野生动物数量更准确的估计，请给出一种你认为更合理的抽样方法，并说明理由.
\end{enumerate}

附：相关系数
\(r=\dfrac{\sum_{i=1}^{n}(x_i-\bar x)(y_i-\bar y)}{\sqrt{\sum_{i=1}^{n}(x_i-\bar x)^2\sum_{i=1}^{n}(y_i-\bar y)^2}}\)，
\(\sqrt2\approx 1.414\).
\end{problem}

\begin{problem}
已知椭圆 \(C_1:\dfrac{x^2}{a^2}+\dfrac{y^2}{b^2}=1\)（\(a>b>0\)）的右焦点 \(F\) 与抛物线 \(C_2\) 的焦点重合，\(C_1\) 的中心与 \(C_2\) 的顶点重合. 过 \(F\) 且与 \(x\) 轴垂直的直线交 \(C_1\) 于 \(A\)，\(B\) 两点，交 \(C_2\) 于 \(C\)，\(D\) 两点，且 \(\lvert CD\rvert=\dfrac{4}{3}\lvert AB\rvert\).

\begin{enumerate}
\item 求 \(C_1\) 的离心率；
\item 若 \(C_1\) 的四个顶点到 \(C_2\) 的准线距离之和为 \(12\)，求 \(C_1\) 与 \(C_2\) 的标准方程.
\end{enumerate}
\end{problem}

\begin{problem}
如图，已知三棱柱 \(ABC-A_1B_1C_1\) 的底面是正三角形，侧面 \(BB_1C_1C\) 是矩形，\(M\)，\(N\) 分别为 \(BC\)，\(B_1C_1\) 的中点，\(P\) 为 \(AM\) 上一点. 过 \(B_1C_1\) 和 \(P\) 的平面交 \(AB\) 于 \(E\)，交 \(AC\) 于 \(F\).

\begin{enumerate}
\item 证明：\(AA_1\parallel MN\)，且平面 \(A_1AMN\perp\) 平面 \(EB_1C_1F\)；
\item 设 \(O\) 为 \(\triangle A_1B_1C_1\) 的中心. 若 \(AO=AB=6\)，\(AO\parallel\) 平面 \(EB_1C_1F\)，且 \(\angle MPN=\dfrac{\pi}{3}\)，求四棱锥 \(B-EB_1C_1F\) 的体积.
\end{enumerate}

\bitmapfigure[float=inline,max width=6.0cm]{img/2020/national_paper_2_liberal/q20_fig1.png}
\end{problem}

\begin{problem}
已知函数 \(f(x)=2\ln x+1\).

\begin{enumerate}
\item 若 \(f(x)\le 2x+c\)，求 \(c\) 的取值范围；
\item 设 \(a>0\)，讨论函数 \(g(x)=\dfrac{f(x)-f(a)}{x-a}\) 的单调性.
\end{enumerate}
\end{problem}

\section{选考题：解答题}

\begin{problem}
已知曲线 \(C_1\) 的参数方程为
\(\begin{cases}
x=4\cos^2\theta,\\
y=4\sin^2\theta
\end{cases}\)（\(\theta\) 为参数），\(C_2\) 的参数方程为
\(\begin{cases}
x=t+\dfrac{1}{t},\\
y=t-\dfrac{1}{t}
\end{cases}\)（\(t\) 为参数）.

\begin{enumerate}
\item 将 \(C_1\)，\(C_2\) 的参数方程化为普通方程；
\item 以坐标原点为极点，\(x\) 轴正半轴为极轴建立极坐标系. 设 \(C_1\)，\(C_2\) 的交点为 \(P\)，求圆心在极轴上，且经过极点和 \(P\) 的圆的极坐标方程.
\end{enumerate}
\end{problem}

\begin{problem}
已知函数 \(f(x)=\lvert x-a^2\rvert+\lvert x-2a+1\rvert\).

\begin{enumerate}
\item 当 \(a=2\) 时，求不等式 \(f(x)\ge 4\) 的解集；
\item 若 \(f(x)\ge 4\)，求 \(a\) 的取值范围.
\end{enumerate}
\end{problem}

